\section{Related Work}
\label{sec:related-work}

We organize related work by conflict points and complementarity with the VSJoin objective: streaming vector similarity joins under multicore sliding-window execution.

\paragraph{(1) Multicore and distributed stream joins for structured data.}
Prior stream processing research provides the execution foundation for correctness and low-latency operation under unbounded inputs~\cite{akidau2015dataflow}. Representative join designs such as Low-latency Handshake Join~\cite{roy2014llhj}, scalable distributed stream join execution~\cite{lin2015scalable}, adaptive distributed stream join models~\cite{wang2024lowlatency}, and PIM-Tree-based multicore joins~\cite{shahvarani2020pimtree} are highly relevant on concurrency control, indexing under updates, and execution-pipeline discipline.

\textit{Conflict point.} Their routing and locality assumptions are primarily key-centric, while vector similarity joins require neighborhood-aware routing and boundary coverage controls.

\textit{Complementarity.} We inherit the stream-execution perspective (window semantics, continuous correctness, low-latency operation), but extend the control surface to vector-specific routing/index maintenance trade-offs.

\paragraph{(2) Similarity join methods and set-similarity streaming lines.}
Approximate vector indexing methods such as product quantization~\cite{jegou2011product} and graph-based search such as HNSW~\cite{malkov2018efficient} provide strong retrieval performance and are essential components in modern vector systems. In parallel, streaming set-similarity efforts study distributed candidate generation and filtering under token/set semantics~\cite{yang2020distributed,rafiei2020similarity,morales2016streaming}. The main conflict with our target is data and maintenance regime: token-set pipelines and static/batch vector assumptions amortize updates differently from continuous insert/expire/probe traffic under sliding windows with dense-vector threshold predicates.

\textit{Complementarity.} VSJoin does not compete with these methods as standalone indexes; instead, it composes index behavior with routing and control-plane policies so quality/latency remain stable under update churn.

\paragraph{(3) Streaming vector pipelines and retrieval-centric workloads.}
LLM-era retrieval workloads increase demand for fresh, low-latency vector processing~\cite{lewis2020rag}. Dynamic vector maintenance systems such as SPFresh~\cite{xu2023spfresh} are close in spirit, but they focus on vector serving and online index freshness at large scale, rather than single-node multicore stream-stream threshold joins with explicit coverage/control co-design.

\textit{Conflict point.} Our focus is not general vector serving or set-overlap filtering; it is pairwise thresholded vector join under window constraints, where boundary fanout, periodic freshness control, and skew-sensitive remapping must be coordinated explicitly.

\paragraph{Summary of distinction.}
Compared with prior work, our contribution is an implementation-aligned co-design of three mechanisms in one runtime: bounded-staleness read/write decoupling, budgeted boundary coverage routing, and predictable control-plane remapping. In our evaluated setting, this combination provides a more direct answer to streaming vector-join tensions than any single prior axis alone (stream runtime, ANN index, or retrieval pipeline).
